\chapter{Conclusion}
\label{cap:conclusion}

This chapter summarizes the work, its contributions, its limitations, and directions for follow-up work.

\section{Summary}
\label{sec:conclusion_summary}

The primary goal of this thesis was to design and implement a real-time speech and music classifier.
The real-time requirement invited strict causality as the self-imposed constraint.
The reference methods address this task as 2-class speech-vs-music, and this thesis extends them to a 3-class formulation (\emph{Speech}, \emph{Music}, \emph{Background}) by adding silence and ambient noise as a third class.

A deterministic dataset construction pipeline was developed and produced a $\sim 100$-hour training dataset from four public sources, paired with a manually labeled critical subset for behavioral review.
Four reference classifiers from the literature were implemented: a decision tree from Lavner and Ruinskiy, a Gaussian mixture model and a support vector machine from Khonglah and Prasanna, and a causal Temporal Convolutional Network (TCN) from Lemaire and Holzapfel.
The causal TCN reached macro F1 score $0.9752$, ten percentage points above the strongest traditional reference.

Two experimental phases followed.
The effectiveness phase explored the TCN's design dimensions and produced TCN-L ($\sim 389$\,K parameters, macro F1 score $0.9846$) as the representative high-capacity variant at roughly $12\times$ the reference parameter count.

After it became clear that further accuracy gains required much larger models, the cost phase reversed direction and asked how much computational cost could be shed at comparable accuracy.
That phase produced \textbf{TCN-S} ($\sim 4.6$\,K parameters, macro F1 score $0.9722$), which matches the reference TCN within $0.003$ macro F1 score at $1/7$ the parameter count and roughly $21\times$ less per-frame compute.
A single-thread CPU benchmark confirmed that every implemented classifier runs above $2\times$ real-time, with TCN-S leading the family at $5.3\times$.

\section{Contributions}
\label{sec:conclusion_contributions}

The primary contribution of this thesis is the proposed \textbf{TCN-S}, a causal classifier ($\sim 4.6$\,K parameters) that matches the reference TCN within $0.003$ macro F1 score at $1/7$ the parameter count and roughly $21\times$ less per-frame compute.
A second contribution is the \textbf{dataset construction pipeline}: a deterministic build that produced the $\sim 100$-hour training dataset from four public sources, designed so that adding new sources or download adapters extends the dataset without redesigning the pipeline.

The remaining contributions support the two above:
\begin{itemize}
    \item \textbf{Implementations of three traditional references} extended from 2-class to 3-class streaming: a decision tree from Lavner and Ruinskiy, and a Gaussian mixture model and a support vector machine from Khonglah and Prasanna.
    \item \textbf{Causal TCN implementation} following Lemaire and Holzapfel, plus the derived TCN-L ($\sim 389$\,K parameters, macro F1 score $0.9846$) as a high-accuracy comparison point.
    \item \textbf{Systematic experimental program} over the TCN's design space: sensitivity ablation, spectral frontend, learned preprocessor, alternative backbones, combined variants, temporal heads, and the cost-reduction sweep that derives TCN-S.
    \item \textbf{Cross-model analyses}: a single-thread CPU benchmark, symbolic complexity metrics (parameters, streaming MACs, receptive field), and a critical-subset walk-through for manual behavioral review.
\end{itemize}

The poster and presentation video required by the assignment are submitted alongside the thesis, together with the source code needed to reproduce or build upon it.

\section{Limitations}
\label{sec:conclusion_limitations}

Two main limitations of this work are noted below.

\textbf{Per-model frame-length differences.}
Each implemented model emits decisions at its own cadence: $10\,\text{ms}$ for the decision tree, $15\,\text{ms}$ for GMM and SVM, and $23.2\,\text{ms}$ for the TCN family.
These hops follow each model's on-line adaptation rather than the reference papers verbatim, and they make the decision tree's predictions look more jittery in side-by-side timeline plots than the others, even when accuracy is comparable.
A uniform-frame-rate re-evaluation (\Cref{sec:conclusion_future}) would remove this artifact.

\textbf{Dataset coverage gap on \emph{Background}.}
The critical-subset review (\Cref{subsec:exp_critical_results}) surfaced a gap in the assembled dataset: ambient sounds the classifier cannot recognize collapse onto \emph{Music}, and TCN-S keeps that label even when speech enters the mix.
The gap traces to the limited \emph{Background} variety in the public sources used and is the main motivation for dataset expansion described in \Cref{sec:conclusion_future}.


\section{Follow-up work}
\label{sec:conclusion_future}

The work leaves several directions open for follow-up. Three are outlined below.

\textbf{Dataset expansion.}
The assembled dataset is large enough to support training but does not yet cover the full variety of inputs the classifier will encounter, with \emph{Background} as the most prominent gap.
Expanding training coverage across all three classes may allow a re-trained classifier to close most of the visible robustness gap.
A separate evaluation dataset drawn from sources disjoint from the training set would in turn quantify generalization beyond the training distribution.

\textbf{Codec integration and perceptual-quality testing.}
The natural next step is integration into a speech or audio codec, e.g., the speech/music classifier of EVS~\cite{dietz2015evs}, where TCN-S sits well within the codec's per-frame CPU budget.
A codec variant that swaps in a classifier from this thesis can then be compared against the production setup by objective and subjective perceptual-quality testing.
This comparison answers the question macro F1 score cannot: whether the proposed classifier yields a better-sounding codec, not just a better-scoring classifier.

\textbf{Uniform frame-rate evaluation.}
The implemented classifiers emit decisions at different cadences (\Cref{sec:conclusion_limitations}).
Re-evaluating every model at one shared frame rate would normalize the visual cadence in side-by-side timeline plots and the per-frame metric weighting, removing the jitter artifact.
